(1) 取质点坐标 \(x=X+s\cos\alpha,\ y=-s\sin\alpha\)。因此 \(T=\frac12M\dot X^2+\frac12m[(\dot X+\dot s\cos\alpha)^2+\dot s^2\sin^2\alpha]\)，\(V=-mg s\sin\alpha\)。它包含速度交叉项 \(m\dot X\dot s\cos\alpha\)，不能把两者动能简单相加为 \(m\dot s^2/2\)。

(2) 坐标 \(X\) 循环，\(P_x=\partial L/\partial\dot X=(M+m)\dot X+m\dot s\cos\alpha\) 恒定。若初态全静止，\(P_x=0\)，故 \(\boxed{\dot X=-m\cos\alpha\,\dot s/(M+m)}\)，楔块向质点相对移动方向的反方向运动。

(3) 对 \(s\) 的Euler–Lagrange方程是 \(m(\ddot s+\ddot X\cos\alpha)=mg\sin\alpha\)。代入 \(\ddot X=-m\cos\alpha\ddot s/(M+m)\) 得 \(\boxed{\ddot s=g\sin\alpha(M+m)/(M+m\sin^2\alpha)}\)，\(\ddot X=-mg\sin\alpha\cos\alpha/(M+m\sin^2\alpha)\)。\(M\to\infty\) 时 \(\ddot s\to g\sin\alpha\)。

(4) 消去 \(\dot X\) 后 \(T=\frac12m(M+m\sin^2\alpha)\dot s^2/(M+m)\)。从静止下降高度 \(h\)，能量守恒给 \(\dot s^2=2gh(M+m)/[M+m\sin^2\alpha]\)。楔的动能占总释放势能份额为 \(T_M/(mgh)=Mm\cos^2\alpha/[(M+m)(M+m\sin^2\alpha)]\)。光滑约束只在两体间传力，总系统没有约束耗散。
